\section{Adapting the Abstract Blockchain Interface}\label{sec:adapting}

\begin{figure*}
  \begin{subfigure}[b]{0.3\textwidth}
    \centering
     \begin{tikzpicture}[scale=1]
      \draw [fill=verylightviolet,thick] (0,0) rectangle +(1.5,1);      
      \draw (0.75,0.5) node {peer};
      \draw[<->] (-1.5,0.5) -- (0,0.5);
      \draw (-0.75,0.7) node {{\scriptsize API calls}};
      \draw (-2.1,0.55) node {clients};
    \end{tikzpicture}
     \caption{A blockchain peer and its clients connected, directly,
       through method calls to the abstract blockchain interface.
       Both peer and clients are running in the same JVM, in the same machine,
       as Java objects: clients can call the methods of the peer.}
     \label{fig:local_peer}
  \end{subfigure}
  \quad
  \begin{subfigure}[b]{0.3\textwidth}
    \begin{tikzpicture}[scale=1]
      \draw [fill=verylightyellow,thick] (-0.2,-0.2) rectangle +(1.9,1.7);
      \draw (0.75,1.25) node {peer service};
      \draw [fill=verylightviolet,thick] (0,0) rectangle +(1.5,1);      
      \draw (0.75,0.5) node {peer};
      \draw[<->] (-2,0.5) -- (-0.2,0.5);
      \draw (-1.1,0.7) node {{\scriptsize network calls}};
      \draw (-2.6,0.55) node {clients};
    \end{tikzpicture}
    \caption{A blockchain peer, published as a network service
      whose abstract blockchain application can be called through network calls,
      whose arguments and results are represented in JSON format.
      Peer and clients can run in different JVM and different machines.}
    \label{fig:web_peer}
  \end{subfigure}
  \quad
  \begin{subfigure}[b]{0.3\textwidth}
    \begin{tikzpicture}[scale=1]
      \draw [fill=verylightred,thick] (-0.4,-0.4) rectangle +(2.3,2.4);
      \draw (0.75,1.75) node {remote peer};
      \draw [fill=verylightyellow,thick] (-0.2,-0.2) rectangle +(1.9,1.7);
      \draw (0.75,1.25) node {peer service};
      \draw [fill=verylightviolet,thick] (0,0) rectangle +(1.5,1);      
      \draw (0.75,0.5) node {peer};
      \draw[<->] (-2.2,0.5) -- (-0.4,0.5);
      \draw (-1.3,0.7) node {{\scriptsize API calls}};
      \draw (-2.8,0.55) node {clients};
    \end{tikzpicture}
    \caption{The network service of a blockchain peer, adapted to the
      abstract blockchain interface as a remote peer. This allows clients
      to interact with the peer with the same method calls used
      in~(a), hiding the fact that the peer
      is remote, potentially running on a different machine.}
    \label{fig:remote_peer}
  \end{subfigure}
  \caption{Adaptation of a peer to a network service, in order to make it accessible
    through network calls from clients running in different machines; its further
    adaptation as a remote peer, in order to use the same
    abstract blockchain interface also for peers running in different machines.}
\end{figure*}

The testcases in \cref{sec:takamaka} can be run as clients against a peer, created
as a Java object in the same JVM where the testcases are executed
(\cref{fig:local_peer}). In this case, the testcases can execute the API
calls in \cref{fig:api} directly. However, another, more realistic, scenario is that
of a peer running in a different machine and published as a web service, that can be
used by clients through network calls corresponding to
the method calls in \cref{fig:api}, whose arguments and results
are represented in JSON format (\cref{fig:web_peer}). This scenario requires to
modify the testcases in \cref{fig:testcases}, in order to replace the direct method
calls in \cref{fig:api} with network calls. It also requires to represent API arguments
in JSON format and back. This modification is problematic since it is against
the open-closed principle: our testcases should be open to extension
(\ie able to run both on a local and a peer reachable thirough a network)
without requiring their
code modification. Concretely, we should maintain two versions of the testcases,
either for testing agains a local peer and another for testing against a peer
accessible through the network.
In order to avoid this software engineering problem, we have adapted
network-reachable peers to the API in \cref{fig:api}, through the adaptor
design pattern~\citep{GammaHJV95}. This is shown in \cref{fig:remote_peer}. A remote peer
is a Java object that implements the API in \cref{fig:api} and adapts
them into network calls over JSON-represented arguments, and back.
A remote peer has the same interface as a local peer, therefore the testcases
in \cref{fig:testcases} needn't be modified to run against a remote peer.

\Cref{sec:experiments} reports experiments with both local and remote peers.
In the latter case, experiments are more realistic since the testcases
cover also the networking layer of the peers and the representation into and from JSON
of the API calls' arguments.
